\documentclass[12pt,a4paper]{ctexrep}

% --- 宏包 ---
\usepackage{geometry,setspace,titlesec,fancyhdr}
\usepackage{amsmath,amssymb,booktabs,tabularx,longtable,graphicx}
\usepackage{enumitem}
\usepackage{fvextra}
\usepackage[hidelinks]{hyperref}
\usepackage{float}
\numberwithin{equation}{section}

% --- 页面边距与布局 ---
\geometry{top=2.5cm,bottom=2.54cm,left=2.5cm,right=2.5cm}

% --- 字体设置 ---
\setCJKmainfont{SimSun}
\setCJKsansfont{SimHei}
\setmainfont{Times New Roman}
\setmonofont{Consolas}
\newcommand{\song}{\songti}
\newcommand{\hei}{\heiti}

% --- 段落与页眉页脚 ---
\setlength{\parskip}{0pt}
% --- 页眉页脚设置 ---
\pagestyle{fancy}
\fancyhf{}
\cfoot{\song\footnotesize\thepage}
\renewcommand{\headrulewidth}{0pt}


% --- 重新定制标题格式 ---
\titleformat{\section}[block]{\hei\zihao{-3}}{\thesection}{1em}{}
\titlespacing{\section}{0pt}{8pt}{8pt}

\titleformat{\subsection}[block]{\hei\zihao{4}}{\thesubsection}{1em}{}
\titlespacing{\subsection}{0pt}{5pt}{5pt}

\titleformat{\subsubsection}[block]{\hei\zihao{4}}{\thesubsubsection}{1em}{}
\titlespacing{\subsubsection}{0pt}{5pt}{5pt}

\renewcommand{\thesection}{\arabic{section}}
\renewcommand{\thesubsection}{\thesection.\arabic{subsection}}
\renewcommand{\thesubsubsection}{\thesubsection.\arabic{subsubsection}}
\setcounter{secnumdepth}{3}

% --- 引用符号定制（上标、带中括号） ---
\newcommand{\citeup}[1]{\textsuperscript{[#1]}}

\begin{document}
	% 前置部分：中文摘要 
	\pagenumbering{gobble}
	\newpage	
	% ===================================================
	% 前置部分：中文摘要 ===================================================
	\newpage
	\pagenumbering{arabic} % 设置阿拉伯数字页码
	
	% --- 1. 中文摘要 ---
	\begin{center}
		{\hei\zihao{2}标题}
	\end{center}
	\noindent{\hei\zihao{-3}摘要}\par
	\vspace{11pt}
	
	\begin{spacing}{1.0}
		{\song\zihao{4} \setlength{\parindent}{2em}
			中文摘要
		}
	\end{spacing}
	
	\vspace{1em}
	\noindent{\hei\zihao{4}关键词：1；2；3；4}
	
	
	%===================================================
	% 正文部分
	%===================================================
	\begin{spacing}{1.10}
		{\song\zihao{-4}
		%引言部分
		\section{引言}
			\subsection{问题背景}
			%简要概括这个研究目标的重要性
			\subsection{问题重述}
			%简要概括问题的具体是什么，转化数学建模目标
			
			\section{问题分析}
			\subsection{数据概况}

			附件由男胎表和女胎表组成，每张表均含31个原始字段，记录孕妇基本信息、检测时间、孕周与BMI、测序质量、染色体统计量和检测结果。男胎表有1082条记录，涉及267位孕妇；女胎表有605条记录，涉及147位孕妇。男、女胎表中分别有19组和15组记录具有相同的候选采血键。由于记录数远多于孕妇数，同一孕妇的多次检测构成重复测量，不能视为相互独立的样本。

			\begin{table}[htbp]
				\centering
				\caption{男胎与女胎检测数据的主要结构}
				\label{tab:data-overview}
				\small
				\renewcommand{\arraystretch}{1.15}
				\begin{tabularx}{\textwidth}{lrrX}
					\toprule
					数据特征 & 男胎 & 女胎 & 分析含义 \\
					\midrule
					检测记录数/孕妇数 & 1082/267 & 605/147 & 需要处理个体内相关并按孕妇划分数据 \\
					孕周范围（天） & 77--203 & 77--199 & 决定模型的观测支持与外推边界 \\
					BMI范围（$\mathrm{kg/m^2}$） & 20.70--46.88 & 25.59--45.72 & 尾部样本较少，不宜继续细分BMI区间 \\
					AB异常记录/涉及孕妇 & --- & 67/44 & 第四问需要兼顾少数类和孕妇隔离 \\
					\bottomrule
				\end{tabularx}
			\end{table}

			男胎中有937条记录达到题设的$4\%$阈值，其余145条低于阈值。从孕妇的观测序列看，218位在首次观测时已经达标，42位在后续观测中首次达标，另有7位在全部观测中均未达标。这些统计量仅反映现有随访记录中的观测状态，不能代替真实的首次达标时间。女胎的67条AB异常记录涉及44位孕妇，其中43位同时存在AB空白和非空记录。第四问因而保留检测级标签，不将同一孕妇固定标记为异常或正常。
			%承接上文进一步分析四个问题的目标、观测机制和模型选择。
			%在此处插入主流程图

			\section{问题假设}
			%问题的重要假设，支持后续建模
			
			\section{符号及标准说明}
			%文章中的重要符号
			
		%建模准备工作
		\newpage
			\section{数据预处理与验证样本构造}
			\subsection{数据与建模的使用关系}

			问题1利用男胎的检测孕周、BMI等指标解释Y染色体浓度，并处理同一孕妇的重复观测。问题2、3在此基础上研究BMI分组与检测时点，还需使用男胎的浓度变化和个体因素。为防止关系模型提前接触下游测试个体，问题1--3采用同一套孕妇划分。问题4则以女胎单次检测时可获得的染色体统计量和质量指标为输入，以检测级AB结果为预测目标；数据划分既要隔离孕妇，也要为少数类保留足够样本。

			数据预处理后，男胎三问共享连续Y染色体浓度、孕周和个体信息；女胎仅使用当次检测可得的指标判定AB结果。孕妇代码只用于描述相关结构和划分数据，不作为预测特征。异常处理也限定在受影响的字段内，不因单列异常而删除整条记录。

			\subsection{异常数据的分层处置}

			本文根据现有信息能否唯一解释异常，将数据问题分为三类，并记录每项处理的依据和信息边界。

			{\hei 可直接处理的数据。} 对能够由题意或同条记录确定的数据直接补算或统一格式。女胎U、V列记为结构性缺失，AB空白依题意填为未检出异常；“周+天”转换为孕天，“$\geq3$”保留为带下界标记的类别，空白字符串记为缺失。数据中唯一一条BMI缺失记录，可由同条记录的身高$162\,\mathrm{cm}$和体重$95\,\mathrm{kg}$确定，按$\mathrm{BMI}=m/(h/100)^2$补算为$36.198750\,\mathrm{kg/m^2}$。这些处理不借用其他孕妇或后续检测的信息；末次月经等无法唯一恢复的字段不作反向填充。

			{\hei 超过参考范围的数据。} 超出题目给出的参考范围并不足以判定录入错误，医学测量中可能存在边界个体和测量波动。因此，尾部BMI、参考区间外GC值、负X染色体浓度及少量非常规孕周均保留原值，不统一删除或缩尾。后续通过覆盖统计和敏感性分析，检验少数极端观测是否主导结论。

			{\hei 存在硬冲突的数据。} 若多个字段相互矛盾且无法判断错误来源，直接改值会引入主观假设。男胎中有71条记录的唯一读段数大于原始读段数，数据中还存在日期顺序、年龄和孕产次冲突。这类字段保留原值并增加异常标签，不做交换、反算或多数投票，也不因此删除整条记录。若模型依赖冲突字段，则根据标签屏蔽相应的建模使用值，或设置敏感性对照；其余字段仍可参与建模。

			\subsection{数据统一与建模接口}

			语义核查后，数据按“原始事实层---建模使用层---质量审计层”组织。原始事实层保留附件数据和记录标识，供结果追溯；建模使用层为10个关键字段设置使用值，以空值表示该字段在当前任务中无法可靠使用；质量审计层记录字段状态、处理原因和信息范围，并单独存放可能涉及后续随访或标签的信息。这种分层使模型避开不可信字段，同时保留同一记录中仍然有效的信息。

			字段格式按统一规则转换：日期采用ISO格式；检测孕周转换为整数孕天，同时保留周数；比例字段以小数存储；孕次“$\geq3$”保留下界及截尾状态；AB派生出二元异常标签以及T13、T18、T21指示量。规范数据包含48个基础与确定性派生字段、47个建模状态字段，共95列。两张表的1687条记录均予保留，未因异常审计删除样本。进入数据划分阶段后，原有95列的160265个单元格保持不变，划分副本仅新增分折编号。

			\subsection{训练集与测试集设计}

			如图\ref{fig:data-structure}所示，同一孕妇的多条检测记录相互关联,{\hei 各条记录之间并不独立}。因此训练集、验证集和测试集均按孕妇划分，测试孕妇及其全部记录在训练阶段不可见。外层分数由此衡量模型对新孕妇的泛化能力，并防止模型借助个体特征或早期检测信息识别测试记录。
			
			\begin{figure}[H]
				\centering
				\includegraphics[width=\textwidth]{../论文图表原稿/ch6_fig1_data_structure.pdf}
				\caption{记录间不独立事实}
				\label{fig:data-structure}
			\end{figure}

			令第$i$条检测记录所属孕妇为$g_i$，孕妇的外层折编号为$s(g_i)\in\{1,\ldots,5\}$，则第$k$轮的测试集和开发集定义为
			\begin{equation}
				\mathcal D_k^{\mathrm{test}}=\{i:s(g_i)=k\},\qquad
				\mathcal D_k^{\mathrm{dev}}=\{i:s(g_i)\neq k\}.
			\end{equation}

			本文采用{\hei 外层5折、内层3折的嵌套设计}。男胎问题1--3共享同一套孕妇名单，每个外层测试折包含53--54位孕妇。女胎数据单独分层，每折包含29--30位孕妇，其中8--9位出现过AB异常记录。填补、标准化、变量与模型选择、阈值校准，以及BMI组界与检测时点估计，均由内层折在外层开发集内完成。确定方案后，模型在全部外层开发集上重新拟合，并仅在外层测试集上评价一次,将模型选择与最终评价分离。异常分布和误差分析也只在划分完成后审计，以保证外层结果反映模型对新孕妇的泛化性能。

			\subsection{建模准备总结}

			图\ref{fig:data-preparation-flow}概括了本章的数据准备流程：先明确男、女胎数据与四个问题的对应关系，再按“直接处理、保留检验、冲突标记”处置异常，随后建立原始事实、建模使用和质量审计三层接口，并按孕妇构造外层5折、内层3折的验证样本。处理后的数据可以追溯原始记录，也能按任务选择字段；所有需要从数据中估计的步骤均限制在训练折内，为后续建模提供一致的验证基础。

			% 流程图定稿设计要点：
			% 1. 采用一条从左到右的四步主线，并直接标出序号：①数据与四问的关系；②三类异常处置；③三层数据接口；④按孕妇嵌套划分。
			% 2. 每个节点同时写清输入、判断或动作、输出，使评委不阅读正文也能复述完整的数据准备逻辑。
			% 3. 第二步内部用三个并列分支，分别写“补算/规范化”“保留/敏感性分析”“原值不动/增加异常标签”。
			% 4. 第三步用三层叠放结构展示原始值、建模值和审计标签；实线进入模型，审计层用虚线并注明“不作特征”。
			% 5. 第四步明确“同一孕妇不跨折”“外5内3”，末端分为男胎问题1--3共享、女胎问题4分层两路。
			% 6. 正式图只保留必要数字和动词，不写长句；颜色与线型均有语义，黑白打印后仍可区分。
			\begin{figure}[H]
				\centering
				\fbox{\begin{minipage}[c][3.5cm][c]{0.91\textwidth}
					\centering
					{\hei 流程图预留位置}\par\vspace{0.6em}
					{\song\small
					（1）数据与四问的关系 $\rightarrow$ （2）三类异常处置\par
					$\rightarrow$ （3）三层数据接口 $\rightarrow$ （4）按孕妇嵌套划分\par\vspace{0.45em}
					定稿图需在节点内直接给出处理动作、输出和防泄漏边界，\par
					使读者无需查阅正文即可理解完整流程。}
				\end{minipage}}
				\caption{数据准备流程}
				\label{fig:data-preparation-flow}
			\end{figure}

		%正式建模	
			\section{问题1建模}
			\label{sec:q1-model}
			\subsection{预实验与建模分析}
			\label{subsec:q1-pilot}

			问题1要求分析男胎Y染色体浓度与检测孕周、BMI等指标的关系，并检验这种关系是否具有统计意义。经过数据清理后，可用于本问的样本为1081条检测记录，来自267位孕妇。多数孕妇接受过多次检测，若直接按记录拟合普通回归，会低估参数的不确定性，也会把孕妇之间的差异误写成同一孕妇随孕周变化的规律。预实验因此围绕两个问题展开：{\hei 数据中是否存在需要建模的孕妇层级，以及孕周、BMI的个体内效应和个体间效应能否合并}。

			预实验沿用外层5折、内层3折的孕妇级划分，全部重复记录随孕妇进入同一折。外层测试折保持封存，结构判断只使用各外层开发集及其内层验证折。依次考察常数基线、含随机截距的线性模型、低自由度样条模型和浅层树模型。五个开发集的空模型组内相关系数为$0.578$--$0.617$，说明Y浓度的变异有相当部分来自孕妇层级，{\hei 逐条记录独立的假设不成立}。内外效应分解后，个体内孕周效应在五个开发集上均为正，范围为$0.293$--$0.336$\,pp/周；个体平均孕周项均为负，范围为$-0.359$--$-0.313$\,pp/周。BMI也出现类似的内外差异：个体间效应稳定为负，而个体内效应的区间跨过0。由此,{\hei 正式模型必须保留孕妇层级，并分别估计个体内和个体间效应}。

			非线性探针没有改变上述建模方向。低自由度样条相对线性结构的预测误差仅下降$1.310\%$，其条件区间上界为$2.043\%$，未达到预设的$3\%$增益门槛。浅层树的点增益为$4.912\%$，但区间下界只有$2.361\%$，固定效应训练--验证误差比为$1.3248$，超过$1.30$的上限。因此，样条模型停止扩容，浅层树只保留作预测和残差诊断；关系解释仍采用低容量的层级模型。

			\subsection{主模型建立}
			\label{subsec:q1-model-establishment}

			设第$i$位孕妇的第$j$次检测孕周为$t_{ij}$，BMI为$b_{ij}$，Y染色体浓度为$Y_{ij}$。为便于解释，$Y_{ij}$由原始比例换算为百分点（percentage point，记为pp）。记该孕妇在现有观测中的平均孕周和平均BMI分别为
			\begin{equation}
				\bar t_i=\frac{1}{n_i}\sum_{j=1}^{n_i}t_{ij},\qquad
				\bar b_i=\frac{1}{n_i}\sum_{j=1}^{n_i}b_{ij},
				\label{eq:q1-person-means}
			\end{equation}
			其中$n_i$是孕妇$i$的有效检测次数。最终采用的P2\_RS\_WB模型写为
			\begin{equation}
			\begin{aligned}
			Y_{ij}={}&\beta_0+\beta_{tb}(\bar t_i-16)
			+\beta_{tw}(t_{ij}-\bar t_i)
			+\beta_{bb}(\bar b_i-32) \\
			&+\beta_{bw}(b_{ij}-\bar b_i)
			+u_{0i}+u_{1i}(t_{ij}-\bar t_i)+\varepsilon_{ij}.
			\end{aligned}
			\label{eq:q1-main-model}
			\end{equation}

			式\eqref{eq:q1-main-model}中，$\beta_0$是在孕妇平均孕周为16周、平均BMI为$32\,\mathrm{kg/m^2}$，且随机效应取0时的平均Y浓度；$\beta_{tb}$和$\beta_{bb}$分别描述孕妇之间平均孕周、平均BMI不同所对应的浓度差异；$\beta_{tw}$和$\beta_{bw}$描述同一孕妇的孕周、BMI偏离其个人均值时，Y浓度随之发生的变化。下标$b$与$w$分别表示between和within。16周是有效记录孕周的中位数，$32\,\mathrm{kg/m^2}$接近样本BMI的中心位置，二者只用于改善截距的可读性，不是临床阈值或BMI分界点。

			随机截距$u_{0i}$允许不同孕妇具有不同的浓度基线，随机斜率$u_{1i}$允许孕周效应因人而异。设
			\begin{equation}
			\begin{pmatrix}u_{0i}\\u_{1i}\end{pmatrix}
			\sim N\!\left[
			\begin{pmatrix}0\\0\end{pmatrix},
			\begin{pmatrix}
			\sigma_0^2 & \rho\sigma_0\sigma_1\\
			\rho\sigma_0\sigma_1 & \sigma_1^2
			\end{pmatrix}
			\right],\qquad
			\varepsilon_{ij}\sim N(0,\sigma_e^2),
			\label{eq:q1-random-effects}
			\end{equation}
			其中$\sigma_0^2$和$\sigma_1^2$分别是孕妇随机截距和孕周随机斜率的方差，$\rho$是二者的相关系数，$\sigma_e^2$是记录层残差方差。该结构同时处理孕妇内相关、孕周效应的个体差异和内外效应混杂。式中的$\bar t_i$与$\bar b_i$依赖完整的回顾性观测序列，只用于解释现有数据，不能作为新孕妇首次检测时的输入。

			\subsection{主模型求解与结果}
			\label{subsec:q1-solution}

			主模型结构通过逐级嵌套比较确定。P0模型只含孕周、BMI共同斜率和孕妇随机截距；P1在P0基础上拆分个体内、个体间效应；P2再加入个体内孕周随机斜率。模型比较阶段采用极大似然估计，以BIC和似然比统计量为主要依据，同时检查五个外层开发集的收敛情况和随机效应协方差。表\ref{tab:q1-model-selection}给出了两次预先规定的比较。

			\begin{table}[H]
				\centering
				\caption{问题1层级模型的逐级比较结果}
				\label{tab:q1-model-selection}
				\small
				\renewcommand{\arraystretch}{1.15}
				\begin{tabularx}{\textwidth}{lXrrX}
					\toprule
					比较 & 新增结构 & $\Delta\mathrm{BIC}$ & 参考LRT的$p$值 & 裁决 \\
					\midrule
					P1相对P0 & 分开孕周、BMI的个体内与个体间效应 & $-61.926$ & $3.30\times10^{-17}$ & 保留内外效应分解 \\
					P2相对P1 & 增加个体内孕周随机斜率 & $-164.147$ & $2.10\times10^{-39}$ & 采用P2\_RS\_WB \\
					\bottomrule
				\end{tabularx}
			\end{table}

			两级比较均支持更完整的结构。P2在全数据和五个开发集上全部收敛，随机效应协方差矩阵的最小特征值为$0.0552$，没有出现奇异拟合。结构冻结后，使用限制性极大似然法在1081条有效记录上估计参数，并以孕妇为重抽样单位完成1000次bootstrap。表\ref{tab:q1-fixed-effects}列出四个预定义效应；经验$p$值采用add-one修正，避免把有限次重抽样中的零尾数写成$p=0$。

			\begin{table}[H]
				\centering
				\caption{P2\_RS\_WB模型的固定效应与聚类bootstrap结果}
				\label{tab:q1-fixed-effects}
				\small
				\renewcommand{\arraystretch}{1.15}
				\begin{tabularx}{\textwidth}{XrrrX}
					\toprule
					效应 & 估计值 & 95\% bootstrap区间 & 经验$p$值 & 统计结论 \\
					\midrule
					个体内孕周（pp/周） & $0.3264$ & $[0.2611,\ 0.3826]$ & $0.001998$ & 显著正相关 \\
					个体平均孕周（pp/周） & $-0.2687$ & $[-0.4165,\ -0.1296]$ & $0.001998$ & 显著；按观察时序解释 \\
					个体平均BMI（pp/$\mathrm{kg/m^2}$） & $-0.2077$ & $[-0.3184,\ -0.0932]$ & $0.001998$ & 显著负相关 \\
					个体内BMI（pp/$\mathrm{kg/m^2}$） & $0.1340$ & $[-0.0736,\ 0.5254]$ & $0.1678$ & 区间跨0，证据不足 \\
					\bottomrule
				\end{tabularx}
			\end{table}

			\begin{figure}[!htbp]
				\centering
				\includegraphics[width=\textwidth]{../论文图表原稿/ch6_fig2_model_mechanism.pdf}
				\caption{P2\_RS\_WB模型的内外效应分解、结构选择与稳健性。\textbf{a} 对称条形图的上半部比较孕周和BMI的组内、组间效应，下半部比较P0给出的共同斜率；条长表示效应绝对值，方向由数值正负号标明。孕周共同斜率为$+0.314$\,pp/周，却掩盖同一孕妇的正向变化与不同孕妇之间的负向差异；BMI共同斜率为$-0.138$\,pp/$\mathrm{kg/m^2}$，容易把孕妇间差异误读为短期个体作用。\textbf{b} 个人均值项$\bar t_i-16$和$\bar b_i-32$估计孕妇间效应，相对个人均值的偏离$t_{ij}-\bar t_i$和$b_{ij}-\bar b_i$估计孕妇内效应；重点结果为个体内孕周效应$+0.326$\,pp/周和个体平均BMI效应$-0.208$\,pp/$\mathrm{kg/m^2}$，个体内BMI效应的区间跨0。\textbf{c} 极大似然估计下，BIC由P0的4920.3降至P1的4858.3，再降至P2的4694.2；相邻模型的$\Delta\mathrm{BIC}$分别为$-61.9$和$-164.1$。\textbf{d} 267位孕妇的1000次整群bootstrap分布、点估计与95\%区间；1000次重抽样全部成功，个体内孕周、个体平均孕周和个体平均BMI三项方向稳定，个体内BMI仍存在不确定性。平均孕周项还包含信息性复测与观察时序影响。}
				\label{fig:q1-model-mechanism}
			\end{figure}

			控制模型中的其余项后，在同一孕妇的现有观测序列内，孕周每增加1周，Y浓度平均增加$0.3264$个百分点；孕妇平均BMI每增加$1\,\mathrm{kg/m^2}$，其平均Y浓度降低$0.2077$个百分点。个体内BMI效应没有通过显著性检验，故不据此判断短期BMI变化会提高或降低Y浓度。模型截距为$7.9726$个百分点；随机截距方差、孕周随机斜率方差和残差方差分别为$6.4335$、$0.06734$和$1.8042$，截距与斜率的相关系数为$0.4132$。这些结果说明孕妇之间的浓度基线和孕周变化速度均有差异。

			个体平均孕周项为负并不表示孕周增加会使Y浓度下降。观察过程模型显示，当前Y浓度每增加1个百分点，之后仍出现更晚检测记录的优势比为$0.7909$，95\%置信区间为$[0.7380,\ 0.8477]$，$p=3.30\times10^{-11}$。复测机会与当前浓度相关，因而$\beta_{tb}$还包含采样时序和随访选择的影响。本问据此将结论限定为同来源、已观察男胎记录中的回顾性关联。

			{\hei 外层泛化评价。} P2结构冻结后，本文启封预先划定的5个外层测试折。每一折只用其余4折重新估计P2和P0，测试孕妇的随机效应统一取0。主指标为孕妇等权RMSE，辅助指标包括MAE、偏差、$R^2$、$Y<4\%$的Brier分数和AUC。相对增益的区间由5000次孕妇配对bootstrap得到。回顾性整段场景使用测试孕妇全部记录的孕周和BMI计算个人均值，但不使用测试Y；首次观测场景每人只保留最早一条记录，并令个人均值等于该次观测、个体内偏差为0。前者检验模型能否概括已经形成的观测序列，后者才接近第二问面对新孕妇时的信息条件。

			\begin{table}[H]
				\centering
				\caption{P2模型的外层折外预测结果}
				\label{tab:q1-generalization}
				\small
				\renewcommand{\arraystretch}{1.15}
				\begin{tabularx}{\textwidth}{Xlrrrrr}
					\toprule
					评价场景 & 模型 & RMSE/pp & 偏差/pp & $R^2$ & Brier & AUC \\
					\midrule
					回顾性整段 & 训练均值 & 3.3344 & -0.0133 & -0.0041 & 0.1130 & 0.4717 \\
					回顾性整段 & P2 & 3.0310 & 0.0216 & 0.1703 & 0.1005 & 0.7690 \\
					\midrule
					首次观测 & 训练均值 & 3.0894 & 1.0539 & -0.1382 & 0.1551 & 0.4715 \\
					首次观测 & P0 & 2.8604 & -0.1147 & 0.0242 & 0.1512 & 0.5704 \\
					首次观测 & P2 & 3.6435 & 2.2537 & -0.5832 & 0.1685 & 0.5904 \\
					\bottomrule
				\end{tabularx}
			\end{table}

			表\ref{tab:q1-generalization}显示，回顾性整段场景下，P2相对训练均值基线将RMSE降低$9.10\%$，95\%区间为$[5.70\%,\ 12.60\%]$，超过预设的$3\%$增益标准。其95\%正态预测区间覆盖率为$96.28\%$。这组结果只支持模型概括同来源的完整观测序列，不能替代首次检测评价。首次观测时，P2的RMSE比训练均值高$17.94\%$，相对改善区间为$[-24.22\%,\ -12.06\%]$，Brier分数也恶化$8.60\%$。P0的首次观测RMSE改善$7.41\%$；Brier改善只有$2.49\%$，且95\%区间$[-5.38\%,\ 8.04\%]$跨过0，尚不能提供稳定的风险概率。

			五折P2回顾性RMSE的最大值与最小值之比为$1.4237$，超过预设上限$1.30$。20周以后的RMSE升至$3.9042$\,pp。超出相应训练折取值范围的样本只有3位孕妇、7条记录，RMSE为$4.2967$\,pp。P2因此保留为第一问的关系模型，回顾性整段预测只作辅助证据，首次观测和低支持区不作预测性外推。本次评价属于选择后的同来源内部折外检验。测试折没有参与各折拟合，但全样本审计和前期结构判断已经提供过整体信息，故不能称为外部验证或完全未查看的盲测。

			\subsection{灵敏性分析与稳健性分析}
			\label{subsec:q1-robustness}

			灵敏性分析首先检查式\eqref{eq:q1-main-model}中的中心值选择。孕周中心取样本最小值、四分位数、中位数、均值和最大值，BMI中心在相同六类位置外加入32，共形成42组组合。42次重拟合全部收敛，非截距系数的最大绝对差为$2.86\times10^{-7}$，条件拟合值与残差的最大差为$7.56\times10^{-6}$个百分点。中心值变化只引起截距按代数关系平移，不改变效应大小、显著性或拟合结果。

			表\ref{tab:q1-robustness}汇总了其余检验。聚类bootstrap完整保留每位孕妇的观测序列；质量调整和候选采血聚合用于检查测序质量与同次采血重复是否改变结论；留一孕妇分析则直接衡量单个受试者的影响。

			\begin{table}[H]
				\centering
				\caption{问题1主模型的灵敏性与稳健性检验}
				\label{tab:q1-robustness}
				\small
				\renewcommand{\arraystretch}{1.15}
				\begin{tabularx}{\textwidth}{lXX}
					\toprule
					检验 & 结果 & 对主结论的影响 \\
					\midrule
					孕妇整群bootstrap & 1000/1000次成功，前三项稳定效应区间均不跨0 & 支持显著性结论 \\
					质量调整与候选采血聚合 & 核心三项均未反向，最大相对变化$2.99\%$ & 未触发预设的$25\%$失稳线 \\
					留一孕妇分析 & 267/267次成功，核心三项均未反向，最大相对变化$13.87\%$ & 结论不由单个孕妇主导 \\
					观察模式与GEE & 方向保持一致；至少4次记录者的个体平均孕周项变化$35.0\%$ & 方向较稳，精确幅度依赖观察过程 \\
					\bottomrule
				\end{tabularx}
			\end{table}

			稳健性检验支持个体内孕周、个体平均孕周和个体平均BMI三个效应的方向，但没有消除模型边界。条件残差的超额峰度为$4.006$，绝对残差与孕周的Spearman相关系数为$0.262$，相邻残差相关系数为$-0.406$，说明正态、同方差和简单序列相关结构并不完全符合数据。本文因此以聚类bootstrap和GEE补充模型标准误，同时保留信息性复测造成的选择偏差。根据当前稳健性分析结果，可以认为内部方向稳健。

			\begin{figure}[H]
				\centering
				\includegraphics[width=\textwidth]{../论文图表原稿/ch6_fig3_model_performance.pdf}
				\caption{P2模型的个体依赖与泛化边界。\textbf{a} 对称雷达图比较回顾性整段与首次观测两种信息条件下，P2相对训练均值的五项指标变化；RMSE、MAE和Brier按相对改善率表示，AUC与$R^2$按百分点变化表示，向外为改善、向内为恶化。指标由5个外层测试折合并后按孕妇等权计算，回顾性整段含267位孕妇、1081条记录，首次观测每位孕妇仅保留最早记录，共267条。\textbf{b} 已知孕妇可由既往检测估计随机截距和随机孕周斜率；完全未见孕妇没有可调用的个体随机效应，只能将其置0并依赖固定效应。\textbf{c} P2条件预测的精确加性SHAP，以各预测分量在1081条记录中的均值为基线；随机截距和随机孕周斜率合计占平均绝对归因的54.8\%。全样本描述性拟合中，加入已知孕妇随机效应后RMSE由2.980降至1.085；该结果不属于折外评价。新孕妇首访折外RMSE为3.644\,pp、偏差为$+2.254$\,pp，说明P2的条件预测优势依赖已知孕妇，不应直接作为首次检测的风险概率接口。}
				\label{fig:q1-performance-boundary}
			\end{figure}

			\subsection{第一问总结}
			\label{subsec:q1-summary}

			第一问建立了包含孕妇随机截距和孕周随机斜率的内外效应分解模型，并用孕妇整群bootstrap检验固定效应。同一孕妇的Y浓度随孕周增加而上升，孕妇平均BMI与Y浓度负相关；个体内BMI效应的区间跨0。中心值扰动、质量调整和逐孕妇影响分析没有改变这些方向。外层评价回答的是另一件事：P2对完整回顾性序列有增益，但首次观测误差较大，折间表现也没有达到预设的稳定性标准。关系结论与预测用途必须分开。

			对第二问而言，可复用的是结构信息，而不是P2给出的首次观测概率。个体内孕周效应、平均BMI的负向关系、随机效应方差和bootstrap区间可用于设定候选曲线、误差情景与敏感性范围；观察过程模型给出的复测依赖则提醒我们，观测到的首次达标不等于真实首次达标。个人平均孕周和平均BMI依赖完整随访，不能放入首次决策模型。P2在首次观测上的RMSE为$3.6435$\,pp，偏差为$2.2537$\,pp，这一结果排除了直接复用P2的做法。

			第二问应在每个外层开发集内另建边际概率模型，只使用检测时已经获得的孕周、BMI和允许的个体信息，输出$P(Y<4\%\mid t,\mathrm{BMI},\boldsymbol{x})$及其误差范围。P0可作为低容量基线，但仍需在内层折完成模型选择和概率校准；BMI组界、检测时点与损失权重也只能在训练范围内确定。最终比较沿用问题1--3共享的孕妇折号，通过折外预测计算Brier分数、校准误差和决策风险。第二问只接收第一问已经验证的关系方向与不确定性，不接收完整随访信息。
			
			\section{结论}
			%文章主结论，最后正面回应题目要求
		
		%模型评价与推广	
			\section{模型推广与评价}
			\subsection{模型优点}
			%正确处理了什么问题，有什么创新行性
			\subsection{模型局限}
			%在什么地方上，由于什么原因，没有达到
			\subsection{模型推广}
			%可以参考什么，哪些应该慎重，
		}
	\end{spacing}
	
	\section{AI使用声明}
	
	% 参考文献部分
	\section{参考文献}	
	\begin{spacing}{1.25} % 1.25 倍行距
		\zihao{5}             % 五号字
	\begin{enumerate}[label={[\arabic*]}, leftmargin=2em, labelwidth=1.5em, labelsep=0.5em, itemsep=0pt, parsep=0pt, partopsep=0pt, topsep=0pt]		
		\item 						
	\end{enumerate}
	\end{spacing}
	
	%附录部分
	%创建appendix文档 
	
	%附录A 目录
	%创建appendix_catelog 附录部分目录
	
	%附录B 直接引用数据
	\input{appendix_data_citation}
	%论文直接引用的重要数据，比如预实验设计和结果
	
	%附录C 重要长表数据
	%appendix_data_adition
	%部分划分结果，比如数据清洗报告等重要但是没有直接引用的数据
	
	%附录D 支撑材料目录
	%appendix_contents 支撑文件目录，给出文件名和用途
	
	%附录E  代码文件
	%appendix_code     代码文件
\end{document}

